\documentclass[12pt, a4paper]{article}
\usepackage[left=0cm,top=0.5cm,right=0cm,bottom=0cm,bindingoffset=0.4cm]{geometry} % margens
\usepackage{tikz} % criação de figuras
\usepackage{subcaption} % posicionamento de figuras
\usepackage{float}
\pagestyle{empty} % remover número de páginas

\begin{document}
\renewcommand\arraystretch{1.5} % altura das linhas da tabela
\begin{table}[h] % tabela para cabeçalho
\begin{tabular}{|p{2cm}|p{10cm}|p{2cm}|p{4.5cm}|}
\hline
\textsf{Disciplina} & & \textsf{Data}  & \\ \hline
\textsf{Docente}   & & \textsf{Página} & \\ \hline
\end{tabular}
\end{table}

\begin{figure}[H] % adiciona ambas as grades e a linha central

\begin{tikzpicture} % grade esquerda
\draw[step=0.5cm,color=black!50] (-1,-1) grid (9,25); % step = altura da linha, color = cor e intensidade por cento, (x,y) = início da grade, grid (w,h) = largura e altura em centímetros
\end{tikzpicture}
\begin{tikzpicture} % linha central
\draw[color=black] (0,-1) -- (0,25); % ponto inicial -- ponto final
\end{tikzpicture}
\begin{tikzpicture} % grade direita
\draw[step=0.5cm,color=black!50] (-1,-1) grid (9,25);
\end{tikzpicture}

\end{figure}

\end{document}